\hypertarget{ux66f4ux65b0ux8bb0ux5f55}{%
\chapter{更新记录}\label{ux66f4ux65b0ux8bb0ux5f55}}

\hypertarget{section}{%
\section{2026-09-02}\label{section}}

\begin{itemize}
\tightlist
\item
  修复 \texttt{31.2\ VLA模型的强化学习} 在 GitHub 上仍存在的数学公式渲染问题：消除被误识别为 Markdown 一级标题的独立等号行，修正紧邻中文标点、星号上标及参考文献链接文字中的未渲染公式；并按 Word 可见原稿补回噪声转移的步长因子，校正 Flow-Noise 内部状态、噪声网络及 Policy Agnostic RL 的参数符号。
\item
  修复 \texttt{32.3\ 状态空间模型}、\texttt{33.1\ 世界-动作模型概述}、\texttt{33.2\ 世界-动作统一生成模型} 和 \texttt{33.3\ 推理时不生成视频的世界-动作模型} 的 GitHub 数学公式渲染问题：消除独立等号造成的 Setext 标题误判，修正行内公式边界、集合定界符及小于号的转义，并按 Word 可见内容逐式核对；同时让 LaTeX 生成器兼容 GitHub 的反引号行内数学语法，并据此清理 \texttt{15.6}、\texttt{21.5} 既有 LaTeX 输出中的反引号。
\item
  补全第31至35章中正文点名成果的来源：\texttt{31.1} 增补 LAPA、UniVLA 与 Moto，\texttt{32.2} 增补 RSSM、RTFM、IC-World、C-SWM、Marble 与 LeWorldModel，\texttt{32.3} 明确 PlaNet/RSSM、Dreamer 3、Dreamer 4 的对应文献，\texttt{33.2} 增补 LingBot-VA、LingBot-VA 2.0、PAR、PhysGen 与 DreamDojo，\texttt{33.4} 增补 CraftNet、Tactile-VLA 与 ForceVLA2，\texttt{34.1} 增补 GigaWorld-Policy，\texttt{35.1} 与 \texttt{35.2} 增补 Jim Fan 在 AI Ascent 2026 的``Robotics' End Game''演讲。
\item
  从用户提供的候选文献中，为 \texttt{34.2\ UMI数据} 选入与正文四项主题直接对应的 MV-UMI（多视角）、ViTaMIn（视触觉）、DexUMI（灵巧手）和机器人模仿学习数据缩放定律研究；未收录与本节现有论述对应关系较弱的候选项。
\item
  同步当前 Word 对 \texttt{32.4} 的更新，将标题改为``联合嵌入预测架构（JEPA）''，并按原稿调整正文首句；保留原有38处图片内容的人工转写结果及11处裁剪对应关系，不重新处理图片，同时更新目录与 LaTeX 版本。
\item
  删除原公开版第4、5部分（原第22至39章），按当前 Word 版本重新建立第4部分``大模型智能体''、第5部分``扩散模型与多模态生成''和第6部分``具身智能与世界模型''，公开范围调整为第22至35章、54节；Word 原稿未作修改。
\item
  逐张检查新第4至第6部分的470个 Word 图片位置及其中142处 DrawingML 裁剪：448处纯文字、公式、表格或说明性截图人工转写为 Markdown/LaTeX，22张确需保留视觉布局的信息作为裁剪后公开图片资源；未使用自动 OCR，也未保留图片重建占位。
\item
  将第3部分 \texttt{19.9\ Delta\ Attention} 仍在引用的 Kimi Delta Attention 架构图迁入第3部分资源目录，避免删除旧第4部分资源后产生断链。
\item
  按 \texttt{第4-6部分参考文献.md} 补入相应参考文献与补充说明，并对新增引用进行逐项核验；同步更新完整目录、学习路径、README、引用元数据、转换配置及 LaTeX 版本。
\item
  新增 \texttt{15.6\ 注意力机制与梯度下降的等价性}，按本地 Word 实际可见内容人工转写文字与公式，并将动态权重更新和固定 Transformer 层的等价关系保留为 1 张公开流程图；补充上下文梯度下降、线性注意力和归纳头相关参考文献。
\item
  删除 \texttt{18.4\ 高维向量的量化（论文）}；将原 \texttt{27.2\ Delta\ Attention} 迁移为 \texttt{19.9}，将原 \texttt{19.9\ 门控注意力（Gated\ Attention）（论文）} 顺延为 \texttt{19.10\ 门控注意力（Gated\ Attention）}，并按当前 Word 标题去掉 \texttt{20.2\ Engram模块}、\texttt{20.4\ Attention\ Residuals} 及 \texttt{21.1} 至 \texttt{21.4} 标题中的``（论文）''，同时修正 \texttt{20.2}、\texttt{20.4} 及 \texttt{21.1} 至 \texttt{21.4} 的公开 Markdown 文件名。
\item
  删除原 \texttt{21.5} 至 \texttt{21.8}，新增合并后的 \texttt{21.5\ 强化微调前沿研究}，按本地 Word 人工转写 MaxRL、ExGRPO 和结合自蒸馏的经验反思强化学习内容，并补充对应三篇参考文献；同步更新目录、公开范围统计与 LaTeX 版本。
\end{itemize}

\hypertarget{section-1}{%
\section{2026-05-22}\label{section-1}}

\begin{itemize}
\tightlist
\item
  同步本地 Word 对 \texttt{16.1\ GPT与预训练} 的文字修改到 Markdown 与 LaTeX，新增 Scaling Law 小节，并补充预训练缩放、Chinchilla、后训练 RL 和推理时计算缩放相关参考文献。
\item
  同步本地 Word 对 \texttt{38.2\ 多模态统一理解-生成模型的输入编码（论文）} 和 \texttt{38.3\ 多模态统一理解-生成模型的主干架构与训练（论文）} 的文字修改到 Markdown 与 LaTeX，补入 SenseNova-U1 相关输入编码、位置编码和训练流程说明，补充 SenseNova-U1 与 DMD 参考文献，并更新 README 总字数统计。
\item
  同步本地 Word 对 \texttt{32.2\ 世界模型的主要技术路线} 的文字修改到 Markdown 与 LaTeX，补入 \texttt{Gemini\ Omni}，并按正文中提到的世界模型成果系统补全已核验参考文献。
\item
  同步本地 Word 对 \texttt{32.1\ 世界模型概况} 和 \texttt{39.4\ 世界-策略统一架构的具身智能（论文）} 的修改到 Markdown 与 LaTeX：\texttt{32.1} 补入物理世界推理说明并转写两张纯文字示例图，\texttt{39.4} 新增 \texttt{Pelican-Unified\ 1.0} 小节并补充对应 arXiv 参考文献。
\item
  将本地 Word 中原 \texttt{29.7\ 自主学习的智能系统展望（论文）} 迁移为 \texttt{32.4\ 结合世界模型的自主学习智能系统展望（论文）}，同步图文修改到 Markdown 与 LaTeX，按 Word 可见图片分类转写文字/公式并保留必要架构图，参考文献保持不变。
\end{itemize}

\hypertarget{section-2}{%
\section{2026-05-13}\label{section-2}}

\begin{itemize}
\tightlist
\item
  同步本地 Word 对 \texttt{32.1\ 世界模型概况} 中世界模型落地方向表述的文字修改到 Markdown 与 LaTeX。
\item
  同步本地 Word 对 \texttt{37.3\ LeWorldModel（论文）} 中动作向量说明、预测器上下文说明和动作数据依赖局限性的文字修改到 Markdown 与 LaTeX，并按 Word 可见裁剪范围校正第二张图片转写内容。
\end{itemize}

\hypertarget{section-3}{%
\section{2026-05-11}\label{section-3}}

\begin{itemize}
\tightlist
\item
  加入The bitter lesson：新增 \texttt{1.6\ The\ Bitter\ Lesson}，同步本地 Word 纯文本内容到 Markdown 和 LaTeX，并更新目录与公开范围统计。
\item
  对已完成图片重建的章节进行 Word 可见裁剪图复核，按裁剪后的实际可见区域校正公开图片资源，避免上传未裁剪原图中的无关内容。
\item
  修正 \texttt{01-05}、\texttt{08-01}、\texttt{14-02}、\texttt{17-06}、\texttt{19-04}、\texttt{26-03}、\texttt{27-01}、\texttt{29-02}、\texttt{29-03}、\texttt{30-01}、\texttt{30-07}、\texttt{36-02}、\texttt{36-05} 中与 Word 原稿不一致的公式、正文表述或推导片段。
\item
  为 \texttt{04-07} 补入 Word 原稿中第二张 HMM 结构图的裁剪后公开图片，并补充相应参考文献。
\item
  将 \texttt{01-02}、\texttt{15-03}、\texttt{16-03}、\texttt{20-01}、\texttt{20-03} 中若干公开图片替换为 Word 实际可见裁剪版本。
\item
  修正 \texttt{38-08} 中 Å 单位的 Markdown 与 LaTeX 写法，避免 LaTeX 编译时出现缺字警告。
\item
  基于更新后的 Markdown 和公开 assets 重新生成 LaTeX 项目内容、图片映射和资源清单。
\end{itemize}

\hypertarget{section-4}{%
\section{2026-05-09}\label{section-4}}

\begin{itemize}
\tightlist
\item
  完成剩余 300 个公开图片重建占位：重建 \texttt{34-04}、\texttt{34-05}、\texttt{34-07}、\texttt{38-02}、\texttt{38-03}、\texttt{38-04}、\texttt{38-05}、\texttt{39-02}、\texttt{39-04}、\texttt{39-05}，按 Word 可见裁剪图转写文字、公式和表格，并仅新增 1 张必要的公开流程图。
\item
  将 \texttt{20-02\ Engram模块} 归为论文阅读笔记，更新文件名、正文标题和目录链接。
\item
  同步本地 Word 对 \texttt{33-01} 和 \texttt{39-01} 的文字/图片调整，并为 \texttt{39-03} 补充已核验参考文献。
\item
  新增 \texttt{33.5\ 具身智能的未来展望}，同步本地 Word 正文并补充已核验参考文献。
\item
  重建 \texttt{34-06}、\texttt{37-01}、\texttt{37-02}、\texttt{37-03}、\texttt{38-06}、\texttt{38-07}、\texttt{38-08}、\texttt{39-01}、\texttt{39-03} 的下一批图片占位，按 Word 可见裁剪图转写文字/公式，并新增 3 张必要的公开视觉图。
\item
  修复 \texttt{29-05}、\texttt{30-05}、\texttt{30-06}、\texttt{30-07}、\texttt{30-08}、\texttt{30-10}、\texttt{31-02}、\texttt{31-03}、\texttt{31-04}、\texttt{33-02}、\texttt{33-04}、\texttt{34-03}、\texttt{37-01} 中若干 GitHub 数学公式渲染问题。
\end{itemize}

\hypertarget{section-5}{%
\section{2026-05-04}\label{section-5}}

\begin{itemize}
\tightlist
\item
  初始化 GitHub 开源版资料库。
\item
  将本地 Word 原稿转换为 Markdown。
\item
  放弃普通 OCR 公式文本，图片位置改为稳定重建占位。
\item
  建立本地图片重建清单，后续从第 1 章开始人工重建图片内容。
\item
  明确图片分类处理策略：视觉图保留为图片资源，纯文字和公式转写为正文。
\item
  上传 \texttt{01-02} 中前 5 张视觉图，并将第 6 张公式图转写为 LaTeX。
\item
  新增项目内 \texttt{ai-note-image-workflow} Skill，用于后续多 agent 分片处理图片重建与 GitHub 推送。
\item
  按当前公开范围同步 Markdown 内容，并移除不再公开的旧章节。
\end{itemize}

\clearpage
